% !TEX program = pdflatex
\documentclass[11pt,a4paper]{article}

\usepackage[margin=2cm]{geometry}
\usepackage[T1]{fontenc}
\usepackage[utf8]{inputenc}
\usepackage[spanish]{babel}
\usepackage{lmodern}
\usepackage{amsmath,amssymb}
\usepackage{mathrsfs}
\usepackage{fancyhdr}
\usepackage[hidelinks]{hyperref}

\setlength{\headheight}{14pt}
\setlength{\headsep}{0.5cm}
\pagestyle{fancy}
\fancyhf{}
\fancyhead[L]{\small Grupo de Estudio: Teoría Analítica de Números}
\fancyhead[R]{\small\today}
\renewcommand{\headrulewidth}{0.4pt}
\fancypagestyle{plain}{%
  \fancyhf{}
  \fancyhead[L]{\small Grupo de Estudio: Teoría Analítica de Números}
  \fancyhead[R]{\small\today}
  \renewcommand{\headrulewidth}{0.4pt}
}

\newcommand{\N}{\mathbb{N}}
\newcommand{\Nzero}{\mathbb{N}_0}

\title{Tarea de la Clase 0}
\author{}
\date{}

\begin{document}
\maketitle
\noindent{\small Autor: Daniel Soto Villada}\par

\begin{enumerate}
    \item Calcule las valoraciones indicadas y escriba la factorización prima de cada entero:
    \begin{enumerate}
        \item $v_2(12)$ y $v_3(12)$;
        \item $v_3(45)$ y $v_5(45)$;
        \item $v_2(56)$ y $v_7(56)$;
        \item $v_p(1)$ para un primo $p$ cualquiera.
    \end{enumerate}
    En cada caso, compruebe que el entero se recupera como el producto de las potencias primas correspondientes.

    \item Para $y=5$, determine el conjunto $\mathbb{P}_5$. Luego indique cuáles de los números
    \[
      1,\;6,\;12,\;14,\;25,\;30,\;49
    \]
    pertenecen a $\mathbb{S}_5$. Justifique cada respuesta usando los factores primos de esos números.

    \item Para $y=5$, escriba $\mathscr{P}_5(s)$ como producto de tres series geométricas. Después, calcule el entero
    $n=2^{k_2}3^{k_3}5^{k_5}$ para cada una de estas ternas y escriba el término correspondiente $n^{-s}$:
    \[
      (k_2,k_3,k_5)=(1,0,0),\qquad (0,1,1),\qquad
      (2,1,0),\qquad (0,0,2).
    \]

    \item Sea $p$ un primo y $k\in\Nzero$. Demuestre, usando la definición de valoración p-ádica, que
    \[
      v_p(p^k)=k.
    \]
    Demuestre también que $v_q(p^k)=0$ si $q$ es un primo distinto de $p$.

    \item Sea $y\ge2$ y asigne a cada primo $p\in\mathbb{P}_y$ un exponente $k_p\in\Nzero$. Demuestre que la regla
    \[
      (k_p)_{p\in\mathbb{P}_y}\longmapsto
      n=\prod_{p\in\mathbb{P}_y}p^{k_p}
    \]
    establece una correspondencia uno a uno entre esas listas de exponentes y los números de $\mathbb{S}_y$.
    Use esta correspondencia y las series geométricas para justificar la igualdad finita
    \[
      \sum_{n\in\mathbb{S}_y}n^{-s}
      =\mathscr{P}_y(s)
      =\prod_{p\in\mathbb{P}_y}(1-p^{-s})^{-1},
      \qquad \Re(s)>1.
    \]
\end{enumerate}

\end{document}